\section{Step 42 Observe a block boundary and a zero-frame call}\label{pcm-wav-s42}
Use the boundary65 recipe and compare frames sixty-three and sixty-four with the table below. Its result needs history from the first sixty-four frames; a reset at the block boundary changes it. Repeat with zero-filter, which supplies a valid empty prefix and no PCM bytes. Then use the direct checking fixture to poison the unused PCM and coefficient addresses; an empty data array alone cannot prove absence of reads.


For the left channel, the retained inputs at frames fifty-seven through sixty-three
are 27699, 29693, 31687, minus 31855, minus 29861, minus 27867 and minus 25873.
The next input is minus 23879. The right-channel values are each 997 greater.
All eight coefficient codes remain 4096. Carry those seven input codes into the
one-sample second block; the already filtered frame sixty-three is not its history.

\tagpdfsetup{table/header-rows={1}}
\begin{center}
  \begin{tabular}{rrr}
    Frame & Output left & Output right \\
    \hline
    63    & $-84$       & 913          \\
    64    & $-6282$     & $-5285$      \\
  \end{tabular}
\end{center}

Use \texttt{s42-1}, then \texttt{s42-2} and stop at \texttt{filter\_return}.

\paragraph{Observe} The sixty-five-frame result is continuous across the block split. The zero-frame call succeeds with zero clips and touches only its prefix and invocation frame.

\paragraph{Check} What distinguishes the block boundary from a new channel?
\begin{itemize}
  \item[A.] Input history continues across blocks of the same channel.
  \item[B.] Both require zero history.
  \item[C.] The coefficient order reverses.
\end{itemize}
